\HMTNarrativeHeading{Action, area, and information}

The Barbero--Immirzi parameter is presented on the basis of the antecedents
of action, orientation, and incidence already constructed. The angular
coordinates and conjugate channels act together with the marked inclusion
between hexadic and octadic configurations. The functional receives these
data and determines its evaluation. The relation between arithmetic and
geometry remains visible in each of its arguments.

The incidence structure organizes which elements are distinguished and how
they are prolonged. A hexad is a support of six positions in the design;
its lift uses a declared dodecad and alignment. The resulting octad belongs
to the binary realization. The area readers preserve the marking that
makes this correspondence meaningful. The figures help to follow the
transport of the sets and their intersections.

Action and clock subsequently make it possible to express the energy
associated with a reading. The Boltzmann constant enters into thermal
transduction; the Barbero parameter and the areal unit enter into the
geometrical increment. The gravitational constant enters the realization
that relates action and Planck area. When these expressions are composed,
the same action section can cancel in the quotient between energy and
area increment. The cancellation exhibits a concrete structural
dependence.

The law of state change also preserves the relative-entropy term and the
reference with respect to which the comparison is made. The information
of a reduced observation depends on the correlations maintained by the
joint state. This distinction extends the argument of the archive: the
content of the whole and that of a partial reader are related through
explicit operations.

Conservation theory thus acquires a definite physical application. The
state, its orientation, the area assigned to it, and the thermal reading
form a composition with identified antecedents. The development preserves
these antecedents and their proofs. The subsequent synthesis will bring
this composition together with the operator balances and with the
calculable example of entanglement.
